Para medir los tiempos de forma justa y poder replicar las pruebas sin problemas, aislamos las
funciones de cada algoritmo usando la librería \texttt{<chrono>} de C++, midiendo todo en
microsegundos para no perder precisión en los tamaños más chicos.

\textbf{Detalles del entorno y hardware:}
\begin{itemize}
\item \textbf{Sistema Operativo:} WSL 2, subsistema Ubuntu en Windows 11.
\item \textbf{Hardware:} Procesador Intel Core i5-10400F con 24 GB de memoria RAM.
\item \textbf{Compilación:} Todo compilado a través de \texttt{Makefile} usando \texttt{g++}
\item \textbf{Procesamiento y gráficos:} Se programaron scripts en Python con \texttt{pandas}, \texttt{matplotlib} y \texttt{seaborn} para crear las entradas de prueba y generar las graficas.
\end{itemize}

\textbf{Casos de prueba analizados:}
\begin{itemize}
    \item \textbf{Ordenamiento de Arreglos):} Probamos tamaños de $N \in \{10, 10^3, 10^5, 10^7\}$ en orden ascendente, descendente y aleatorio, combinados con los dominios $D1$ y $D7$.
    \item \textbf{Matrices:} Probamos dimensiones de $N \in \{16, 64, 256, 1024\}$ para matrices densas, dispersas y diagonales, usando valores en los dominios $D0$ y $D10$.
\end{itemize}
